\section{Background and hypotheses}
\label{sec:background}

\subsection{State-space formulation}
\label{sec:statespace}

Let $\xv \in \mathbb{R}^d$ denote the spatiotemporal physical state and
$\zv \in \mathbb{R}^m$ an external forcing. Following the flow-matching
convention we write $\xv_1$ for the \emph{target} state, reserving $\xv_0$ for a
base sample and $\xv_\tau$ for the interpolant introduced in
Section~\ref{sec:methods}. The discrete-time dynamics are
\begin{equation}
  \label{eq:dynamics}
  \xv_{1,k+1} \mid \zv \;=\; \Mop_{\Delta}\!\left(\xv_{1,k}, \zv_{k}, \thv\right) + \xi_k ,
  \qquad k \in \{0,\dots,K\},
\end{equation}
with $\Mop_{\Delta}$ the time-stepping operator over a physical step $\Delta$,
$\thv$ the physical parameters, and $\xi_k$ a state-dependent noise process.
Observations are gathered from the state alone,
\begin{equation}
  \label{eq:obs}
  \yv_k = \Hop(\xv_{1,k}) + \epsv_k , \qquad \epsv_k \sim \mathcal{N}(0, \mathbf{R}).
\end{equation}

Writing $\Phi(\xv_1, \zv; \thv)(t_k) = \Mop_\Delta(\xv_1(t_{k-1}), \zv(t_{k-1}); \thv)$
for $k>0$ and $\Phi(\cdot)(t_0) = \xv_1(t_0)$, the Bayesian state-space model is
\begin{equation}
  \label{eq:ssm}
  \left\{
  \begin{array}{rcl}
    \zv &\sim& \pi_{\zv}, \qquad \thv \;\sim\; \pi_{\thv} \\[0.35em]
    \xv_1 &\sim& \pi_{\xv_1}, \quad
      \pi_{\xv_1}\!\left(\xv_1 \mid \zv, \thv\right) \propto
      \exp\!\left[-\tfrac{1}{2}\left\|\xv_1 - \Phi(\xv_1, \zv; \thv)\right\|^2_{\mathbf{Q}}\right] \\[0.35em]
    \yv &=& \Hop(\xv_1) + \epsv_{\mathrm{obs}}, \quad \epsv_{\mathrm{obs}} \sim \mathcal{N}(0,\mathbf{R}),
  \end{array}\right.
\end{equation}
where $\mathbf{Q}$ is the model-error covariance. The prior on the state is thus
\emph{generated by the model}: $\pi_{\xv_1}$ is defined through $\Phi$, and there
is no other route by which prior information enters. This is the formal content
of H2 and H3b below.

\paragraph{Context, and what this paper estimates.}
In operational practice neither the forcing nor the parameters are known
exactly; the system has access to corrupted surrogates $\zv^{*}$ and $\thv^{*}$.
We collect everything the estimator may condition on into the \emph{context}
\begin{equation}
  \label{eq:context}
  \Cv = \left(\yv_{1:K},\; \zv^{*},\; \thv^{*}\right),
  \qquad
  \Cv_{\leq k} = \left(\yv_{1:k},\; \zv^{*},\; \thv^{*}\right).
\end{equation}
\textbf{Throughout this paper the inference target is the state alone}: we study
the posterior $p(\xv_1 \mid \Cv)$ and its filtered counterpart
$p(\xv_1 \mid \Cv_{\leq k})$. Joint state--parameter inference --- sampling
$p(\xv_1, \thv \mid \Cv)$ --- is a different problem with a different cost, and
is out of scope. Where $\thv^{*}$ appears it is an \emph{input} the estimator may
use, never a quantity it estimates.

\subsection{Four structural hypotheses}
\label{sec:hypotheses}

Classical DA rests on commitments rarely written down as hypotheses.

\begin{description}[leftmargin=1.5em,style=nextline]
  \item[\textbf{H1 (Markovianity).}]
    The joint inference problem over a window --- or a whole reanalysis --- is
    replaced by a \emph{sequence of simpler problems}, each conditionally
    independent of the past given the current state. When $\zv^{*}=\zv$ and
    $\thv^{*}=\thv$, the posterior factorises,
    \begin{equation}
      \label{eq:markov}
      p\!\left(\xv_{1,k'\geq k} \mid \yv_{1:k}, \zv, \thv\right)
      = p\!\left(\xv_{1,k'>k} \mid \xv_{1,k}, \zv, \thv\right)\cdot
        p\!\left(\xv_{1,k} \mid \yv_{1:k}, \zv, \thv\right),
    \end{equation}
    and cycling is exact. This is what licenses filtering; 4D-Var and smoothers
    widen the unit to one window but keep the recursion between windows.
    \emph{The factorisation~\eqref{eq:markov} holds only under those idealised
    conditions.} With $\zv^{*}\neq\zv$ or $\thv^{*}\neq\thv$ the conditional
    independence collapses and the memory induced by the corrupted forcing makes
    the process non-Markovian.
  \item[\textbf{H2 (model fidelity is the route to skill).}]
    The better the dynamical model, the better the assimilation.
  \item[\textbf{H3a (ODE/PDE state representation).}]
    The state is a point in the phase space of~\eqref{eq:dynamics}; uncertainty
    is a perturbation of that point --- a covariance, or an ensemble.
  \item[\textbf{H3b (autoregressive error propagation).}]
    Because $\pi_{\xv_1}$ is generated by integrating $\Phi$, every prior error
    is carried forward by the same recursion that carries the signal, and
    compounds at the system's Lyapunov rate.
\end{description}

\begin{table}[t]
\centering
\small
\begin{tabular}{@{}p{0.06\textwidth}p{0.26\textwidth}p{0.24\textwidth}p{0.32\textwidth}@{}}
\toprule
& \textbf{asserts} & \textbf{buys} & \textbf{costs} \\
\midrule
\textbf{H1} & the joint problem factorises as in~\eqref{eq:markov} & $O(1)$ memory,
operational streaming, no training set; robustness to a change of observing
system without retraining & cannot learn from the \emph{observation
configuration}; no amortisation across windows; the factorisation fails exactly
when $\zv^{*}\neq\zv$ \\[0.4em]
\textbf{H2} & skill is bounded by, and increasing in, model fidelity & a
physically consistent, transferable prior & the model \emph{is} the prior, so
model error enters with no buffer \\[0.4em]
\textbf{H3a} & state $=$ point in phase space; uncertainty $=$ perturbation &
tractable propagation; interpretability & the choice of state variable fixes the
conditioning of the inverse problem; a Gaussian update implies
$\PsiNG \equiv 0$ \\[0.4em]
\textbf{H3b} & prior errors propagate by the same integration as the signal &
one mechanism for signal and uncertainty alike & error compounds at the Lyapunov
rate, and \emph{model} sensitivity is inherited by the analysis whether or not
the \emph{posterior} has it \\
\bottomrule
\end{tabular}
\caption{The four hypotheses, what each buys, and what each costs.}
\label{tab:hypotheses}
\end{table}

\paragraph{H2 is load-bearing, and not simply false.}
Better models do assimilate better. The claim is that H2 is \emph{binding}:
because the model supplies the entire prior via~\eqref{eq:ssm}, DA has no
mechanism to absorb model discrepancy. Model-error covariances, inflation and
weak-constraint formulations exist precisely because this is understood; the
claim is that the available buffers are weak, which
Sections~\ref{sec:inert}--\ref{sec:marginal} measure rather than assert.

\paragraph{H3b is what makes H2 bind.}
The two are not independent. Section~\ref{sec:sensitivity} shows that the
apparent force of ``better model $\Rightarrow$ better DA'' is partly an artefact
of the representation: it is the autoregressive propagation that routes all prior
information through $\Phi$, and so makes the analysis inherit the model's
Lyapunov-amplified sensitivity whether or not the posterior has it.
